\clearpage\phantomsection\label{couverture-programme}\pdfbookmark[0]{Le programme dans le livre}{bm-couverture-programme}\gdef\otChapNum{0}\pagestyle{otFront}\otPageTitre{Le programme dans le livre}{Quatrième · Nouveau programme 2026 · Lecture par notions}
Cette annexe relie le programme officiel au livre. Les pages suivantes
associent d’abord les quatre repères de chaque chapitre à des exercices
précis, puis chaque objectif d’apprentissage de l’annexe 2 du programme
aux chapitres qui le travaillent. Les séries indiquées ne sont pas
exhaustives ; les cours et les défis complètent ce parcours.\par\bigskip
Points à retrouver : multiplication des relatifs justifiée par la distributivité (1),
inverse et fraction d’une fraction (2), produits de puissances (3), encadrement
d’une racine (4), preuves littérales (5), conjectures au tableur (6), ratios et
partages (7), coefficient multiplicateur (8), programmes remontés et notation
$P(t)$ (9), étendue et valeur extrême (10), deux épreuves et fluctuation (11),
réciproque et contraposée (12), translations (13), droite des milieux (14),
Pythagore (15), cercle circonscrit (16), pyramides et cônes (17), conditions
et variable (18).\par\bigskip
\begin{attention}Les exercices à trois pastilles et les défis sont des choix d’entraînement. Ils ne changent pas le programme officiel et ne constituent pas une liste d’exigences supplémentaires.\end{attention}
\clearpage
\otPageTitre{Repères et exercices}{Chapitres 1 à 3}\fontsize{10}{13}\selectfont
\otIntertitre{1. Multiplier et diviser des nombres relatifs}\begin{tabularx}{\linewidth}{@{}Yl@{}}\toprule
Produit de deux relatifs et règle des signes & \hyperref[exo-1-1]{1.1}, \hyperref[exo-1-2]{1.2}, \hyperref[exo-1-4]{1.4}, \hyperref[exo-1-6]{1.6}, \hyperref[exo-1-14]{1.14}\\[2mm]
Quotients et multiplications à trou & \hyperref[exo-1-3]{1.3}, \hyperref[exo-1-9]{1.9}, \hyperref[exo-1-13]{1.13}\\[2mm]
Priorités et enchaînements d’opérations & \hyperref[exo-1-7]{1.7}, \hyperref[exo-1-8]{1.8}, \hyperref[exo-1-10]{1.10}, \hyperref[exo-1-15]{1.15}\\[2mm]
Vocabulaire et programmes de calcul & \hyperref[exo-1-5]{1.5}, \hyperref[exo-1-11]{1.11}, \hyperref[exo-1-12]{1.12}\\[2mm]
\bottomrule\end{tabularx}\par\vspace{4mm}
\otIntertitre{2. Nombres rationnels et fractions}\begin{tabularx}{\linewidth}{@{}Yl@{}}\toprule
Simplifier et écrire un rationnel & \hyperref[exo-2-1]{2.1}, \hyperref[exo-2-2]{2.2}, \hyperref[exo-2-4]{2.4}\\[2mm]
Opposé, inverse et quotients & \hyperref[exo-2-3]{2.3}, \hyperref[exo-2-7]{2.7}, \hyperref[exo-2-8]{2.8}, \hyperref[exo-2-14]{2.14}\\[2mm]
Produits et fractions d’une quantité & \hyperref[exo-2-5]{2.5}, \hyperref[exo-2-6]{2.6}, \hyperref[exo-2-12]{2.12}, \hyperref[exo-2-13]{2.13}\\[2mm]
Enchaînements et problèmes de partage & \hyperref[exo-2-9]{2.9}, \hyperref[exo-2-10]{2.10}, \hyperref[exo-2-11]{2.11}, \hyperref[exo-2-15]{2.15}\\[2mm]
\bottomrule\end{tabularx}\par\vspace{4mm}
\otIntertitre{3. Puissances d’exposant positif}\begin{tabularx}{\linewidth}{@{}Yl@{}}\toprule
Définition et calcul d’une puissance & \hyperref[exo-3-1]{3.1}, \hyperref[exo-3-2]{3.2}, \hyperref[exo-3-4]{3.4}\\[2mm]
Puissances de dix et préfixes & \hyperref[exo-3-3]{3.3}, \hyperref[exo-3-5]{3.5}, \hyperref[exo-3-13]{3.13}\\[2mm]
Produits de puissances & \hyperref[exo-3-6]{3.6}, \hyperref[exo-3-7]{3.7}, \hyperref[exo-3-8]{3.8}, \hyperref[exo-3-9]{3.9}, \hyperref[exo-3-10]{3.10}\\[2mm]
Problèmes de répétition et grands nombres & \hyperref[exo-3-11]{3.11}, \hyperref[exo-3-12]{3.12}, \hyperref[exo-3-14]{3.14}, \hyperref[exo-3-15]{3.15}\\[2mm]
\bottomrule\end{tabularx}\par\vspace{4mm}
\clearpage
\otPageTitre{Repères et exercices}{Chapitres 4 à 6}\fontsize{10}{13}\selectfont
\otIntertitre{4. Carrés et racines carrées}\begin{tabularx}{\linewidth}{@{}Yl@{}}\toprule
Carrés parfaits de 0 à 144 & \hyperref[exo-4-1]{4.1}, \hyperref[exo-4-2]{4.2}\\[2mm]
Définition de la racine carrée & \hyperref[exo-4-3]{4.3}, \hyperref[exo-4-5]{4.5}, \hyperref[exo-4-9]{4.9}\\[2mm]
Encadrements et valeurs approchées & \hyperref[exo-4-4]{4.4}, \hyperref[exo-4-6]{4.6}, \hyperref[exo-4-8]{4.8}, \hyperref[exo-4-10]{4.10}, \hyperref[exo-4-14]{4.14}\\[2mm]
Aires et longueurs exactes & \hyperref[exo-4-7]{4.7}, \hyperref[exo-4-11]{4.11}, \hyperref[exo-4-12]{4.12}, \hyperref[exo-4-13]{4.13}, \hyperref[exo-4-15]{4.15}\\[2mm]
\bottomrule\end{tabularx}\par\vspace{4mm}
\otIntertitre{5. Expressions, distributivité et preuves}\begin{tabularx}{\linewidth}{@{}Yl@{}}\toprule
Structure, substitution et réduction & \hyperref[exo-5-1]{5.1}, \hyperref[exo-5-2]{5.2}, \hyperref[exo-5-3]{5.3}, \hyperref[exo-5-10]{5.10}\\[2mm]
Développer par distributivité simple & \hyperref[exo-5-4]{5.4}, \hyperref[exo-5-6]{5.6}, \hyperref[exo-5-7]{5.7}, \hyperref[exo-5-11]{5.11}\\[2mm]
Factoriser avec un facteur commun & \hyperref[exo-5-5]{5.5}, \hyperref[exo-5-8]{5.8}, \hyperref[exo-5-13]{5.13}\\[2mm]
Formules, programmes et preuves & \hyperref[exo-5-9]{5.9}, \hyperref[exo-5-12]{5.12}, \hyperref[exo-5-14]{5.14}, \hyperref[exo-5-15]{5.15}\\[2mm]
\bottomrule\end{tabularx}\par\vspace{4mm}
\otIntertitre{6. Équations du premier degré}\begin{tabularx}{\linewidth}{@{}Yl@{}}\toprule
Inconnue et test d’une solution & \hyperref[exo-6-1]{6.1}, \hyperref[exo-6-9]{6.9}\\[2mm]
Équations du premier degré & \hyperref[exo-6-2]{6.2}, \hyperref[exo-6-3]{6.3}, \hyperref[exo-6-4]{6.4}, \hyperref[exo-6-6]{6.6}, \hyperref[exo-6-7]{6.7}, \hyperref[exo-6-8]{6.8}\\[2mm]
Mise en équation et vérification & \hyperref[exo-6-5]{6.5}, \hyperref[exo-6-11]{6.11}, \hyperref[exo-6-12]{6.12}, \hyperref[exo-6-13]{6.13}, \hyperref[exo-6-14]{6.14}\\[2mm]
Tableur, conjecture et interprétation & \hyperref[exo-6-10]{6.10}, \hyperref[exo-6-15]{6.15}\\[2mm]
\bottomrule\end{tabularx}\par\vspace{4mm}
\clearpage
\otPageTitre{Repères et exercices}{Chapitres 7 à 9}\fontsize{10}{13}\selectfont
\otIntertitre{7. Proportionnalité, rapports et ratios}\begin{tabularx}{\linewidth}{@{}Yl@{}}\toprule
Quatrième proportionnelle et produits en croix & \hyperref[exo-7-1]{7.1}, \hyperref[exo-7-2]{7.2}, \hyperref[exo-7-3]{7.3}, \hyperref[exo-7-6]{7.6}, \hyperref[exo-7-11]{7.11}\\[2mm]
Rapports, ratios et échelles & \hyperref[exo-7-4]{7.4}, \hyperref[exo-7-5]{7.5}, \hyperref[exo-7-15]{7.15}\\[2mm]
Partages proportionnels & \hyperref[exo-7-9]{7.9}, \hyperref[exo-7-10]{7.10}, \hyperref[exo-7-12]{7.12}, \hyperref[exo-7-14]{7.14}\\[2mm]
Graphiques et grandeurs non proportionnelles & \hyperref[exo-7-7]{7.7}, \hyperref[exo-7-8]{7.8}, \hyperref[exo-7-13]{7.13}\\[2mm]
\bottomrule\end{tabularx}\par\vspace{4mm}
\otIntertitre{8. Pourcentages et grandeurs quotients}\begin{tabularx}{\linewidth}{@{}Yl@{}}\toprule
Pourcentages et coefficient multiplicateur & \hyperref[exo-8-1]{8.1}, \hyperref[exo-8-2]{8.2}, \hyperref[exo-8-3]{8.3}, \hyperref[exo-8-6]{8.6}, \hyperref[exo-8-7]{8.7}, \hyperref[exo-8-13]{8.13}\\[2mm]
Évolutions successives et réciproques & \hyperref[exo-8-11]{8.11}, \hyperref[exo-8-12]{8.12}\\[2mm]
Vitesses, durées et distances & \hyperref[exo-8-4]{8.4}, \hyperref[exo-8-5]{8.5}, \hyperref[exo-8-8]{8.8}, \hyperref[exo-8-14]{8.14}\\[2mm]
Débits, masses volumiques et conversions & \hyperref[exo-8-9]{8.9}, \hyperref[exo-8-10]{8.10}, \hyperref[exo-8-15]{8.15}\\[2mm]
\bottomrule\end{tabularx}\par\vspace{4mm}
\otIntertitre{9. Programmes de calcul et dépendance}\begin{tabularx}{\linewidth}{@{}Yl@{}}\toprule
Programmes de calcul et expressions & \hyperref[exo-9-1]{9.1}, \hyperref[exo-9-2]{9.2}, \hyperref[exo-9-11]{9.11}\\[2mm]
Remonter un programme & \hyperref[exo-9-3]{9.3}, \hyperref[exo-9-11]{9.11}\\[2mm]
Formules, notation $P(t)$ et tableaux & \hyperref[exo-9-4]{9.4}, \hyperref[exo-9-10]{9.10}, \hyperref[exo-9-12]{9.12}, \hyperref[exo-9-13]{9.13}, \hyperref[exo-9-15]{9.15}\\[2mm]
Lectures graphiques et proportionnalité & \hyperref[exo-9-5]{9.5}, \hyperref[exo-9-6]{9.6}, \hyperref[exo-9-7]{9.7}, \hyperref[exo-9-8]{9.8}, \hyperref[exo-9-9]{9.9}, \hyperref[exo-9-14]{9.14}\\[2mm]
\bottomrule\end{tabularx}\par\vspace{4mm}
\clearpage
\otPageTitre{Repères et exercices}{Chapitres 10 à 12}\fontsize{10}{13}\selectfont
\otIntertitre{10. Moyenne, médiane et étendue}\begin{tabularx}{\linewidth}{@{}Yl@{}}\toprule
Effectifs, fréquences et moyenne pondérée & \hyperref[exo-10-1]{10.1}, \hyperref[exo-10-2]{10.2}, \hyperref[exo-10-5]{10.5}, \hyperref[exo-10-6]{10.6}, \hyperref[exo-10-11]{10.11}\\[2mm]
Médiane d’une série de petit effectif & \hyperref[exo-10-3]{10.3}, \hyperref[exo-10-7]{10.7}, \hyperref[exo-10-8]{10.8}, \hyperref[exo-10-12]{10.12}, \hyperref[exo-10-15]{10.15}\\[2mm]
Étendue et valeurs extrêmes & \hyperref[exo-10-4]{10.4}, \hyperref[exo-10-9]{10.9}, \hyperref[exo-10-13]{10.13}\\[2mm]
Comparaison de séries et tableur & \hyperref[exo-10-10]{10.10}, \hyperref[exo-10-14]{10.14}\\[2mm]
\bottomrule\end{tabularx}\par\vspace{4mm}
\otIntertitre{11. Probabilités : événements et deux épreuves}\begin{tabularx}{\linewidth}{@{}Yl@{}}\toprule
Univers, événements et notations & \hyperref[exo-11-1]{11.1}, \hyperref[exo-11-3]{11.3}, \hyperref[exo-11-4]{11.4}, \hyperref[exo-11-10]{11.10}, \hyperref[exo-11-14]{11.14}\\[2mm]
Contraire, réunion et intersection & \hyperref[exo-11-2]{11.2}, \hyperref[exo-11-5]{11.5}, \hyperref[exo-11-6]{11.6}, \hyperref[exo-11-15]{11.15}\\[2mm]
Expériences à deux épreuves & \hyperref[exo-11-7]{11.7}, \hyperref[exo-11-8]{11.8}, \hyperref[exo-11-9]{11.9}\\[2mm]
Fréquences et fluctuation & \hyperref[exo-11-11]{11.11}, \hyperref[exo-11-12]{11.12}, \hyperref[exo-11-13]{11.13}\\[2mm]
\bottomrule\end{tabularx}\par\vspace{4mm}
\otIntertitre{12. Droites remarquables et raisonnement}\begin{tabularx}{\linewidth}{@{}Yl@{}}\toprule
Droites remarquables d’un triangle & \hyperref[exo-12-1]{12.1}, \hyperref[exo-12-5]{12.5}, \hyperref[exo-12-10]{12.10}, \hyperref[exo-12-11]{12.11}\\[2mm]
Hypothèse, conclusion et réciproque & \hyperref[exo-12-2]{12.2}, \hyperref[exo-12-4]{12.4}, \hyperref[exo-12-9]{12.9}\\[2mm]
Contraposée, contre-exemple et absurde & \hyperref[exo-12-6]{12.6}, \hyperref[exo-12-8]{12.8}, \hyperref[exo-12-12]{12.12}, \hyperref[exo-12-14]{12.14}\\[2mm]
Preuves avec des aires et des distances & \hyperref[exo-12-3]{12.3}, \hyperref[exo-12-7]{12.7}, \hyperref[exo-12-13]{12.13}, \hyperref[exo-12-15]{12.15}\\[2mm]
\bottomrule\end{tabularx}\par\vspace{4mm}
\clearpage
\otPageTitre{Repères et exercices}{Chapitres 13 à 15}\fontsize{10}{13}\selectfont
\otIntertitre{13. Translations et parallélogrammes}\begin{tabularx}{\linewidth}{@{}Yl@{}}\toprule
Construire l’image d’une figure & \hyperref[exo-13-1]{13.1}, \hyperref[exo-13-2]{13.2}, \hyperref[exo-13-3]{13.3}, \hyperref[exo-13-10]{13.10}, \hyperref[exo-13-11]{13.11}\\[2mm]
Conservations, droites et angles & \hyperref[exo-13-4]{13.4}, \hyperref[exo-13-9]{13.9}, \hyperref[exo-13-14]{13.14}, \hyperref[exo-13-15]{13.15}\\[2mm]
Translations et parallélogrammes & \hyperref[exo-13-6]{13.6}, \hyperref[exo-13-7]{13.7}, \hyperref[exo-13-12]{13.12}\\[2mm]
Reconnaître transformations, frises et pavages & \hyperref[exo-13-5]{13.5}, \hyperref[exo-13-8]{13.8}, \hyperref[exo-13-13]{13.13}\\[2mm]
\bottomrule\end{tabularx}\par\vspace{4mm}
\otIntertitre{14. La droite des milieux}\begin{tabularx}{\linewidth}{@{}Yl@{}}\toprule
Droite des milieux et parallélisme & \hyperref[exo-14-2]{14.2}, \hyperref[exo-14-5]{14.5}, \hyperref[exo-14-8]{14.8}, \hyperref[exo-14-10]{14.10}\\[2mm]
Segment des milieux et longueurs & \hyperref[exo-14-1]{14.1}, \hyperref[exo-14-3]{14.3}, \hyperref[exo-14-6]{14.6}, \hyperref[exo-14-9]{14.9}\\[2mm]
Parallèle passant par un milieu & \hyperref[exo-14-4]{14.4}, \hyperref[exo-14-7]{14.7}, \hyperref[exo-14-12]{14.12}\\[2mm]
Démonstrations en plusieurs étapes & \hyperref[exo-14-11]{14.11}, \hyperref[exo-14-13]{14.13}, \hyperref[exo-14-14]{14.14}, \hyperref[exo-14-15]{14.15}\\[2mm]
\bottomrule\end{tabularx}\par\vspace{4mm}
\otIntertitre{15. Le théorème de Pythagore}\begin{tabularx}{\linewidth}{@{}Yl@{}}\toprule
Calculer une hypoténuse & \hyperref[exo-15-1]{15.1}, \hyperref[exo-15-3]{15.3}, \hyperref[exo-15-7]{15.7}, \hyperref[exo-15-9]{15.9}\\[2mm]
Calculer un côté de l’angle droit & \hyperref[exo-15-2]{15.2}, \hyperref[exo-15-4]{15.4}, \hyperref[exo-15-8]{15.8}, \hyperref[exo-15-12]{15.12}\\[2mm]
Réciproque et contraposée & \hyperref[exo-15-5]{15.5}, \hyperref[exo-15-6]{15.6}, \hyperref[exo-15-10]{15.10}\\[2mm]
Applications et précision des résultats & \hyperref[exo-15-11]{15.11}, \hyperref[exo-15-13]{15.13}, \hyperref[exo-15-14]{15.14}, \hyperref[exo-15-15]{15.15}\\[2mm]
\bottomrule\end{tabularx}\par\vspace{4mm}
\clearpage
\otPageTitre{Repères et exercices}{Chapitres 16 à 18}\fontsize{10}{13}\selectfont
\otIntertitre{16. Triangle rectangle et cercle circonscrit}\begin{tabularx}{\linewidth}{@{}Yl@{}}\toprule
Centre du cercle et milieu de l’hypoténuse & \hyperref[exo-16-1]{16.1}, \hyperref[exo-16-4]{16.4}, \hyperref[exo-16-13]{16.13}\\[2mm]
Médiane et moitié de l’hypoténuse & \hyperref[exo-16-2]{16.2}, \hyperref[exo-16-6]{16.6}\\[2mm]
Réciproque : triangle inscrit dans un demi-cercle & \hyperref[exo-16-3]{16.3}, \hyperref[exo-16-5]{16.5}, \hyperref[exo-16-7]{16.7}, \hyperref[exo-16-9]{16.9}, \hyperref[exo-16-10]{16.10}, \hyperref[exo-16-15]{16.15}\\[2mm]
Constructions sans équerre et problèmes & \hyperref[exo-16-8]{16.8}, \hyperref[exo-16-11]{16.11}, \hyperref[exo-16-12]{16.12}, \hyperref[exo-16-14]{16.14}\\[2mm]
\bottomrule\end{tabularx}\par\vspace{4mm}
\otIntertitre{17. Pyramides et cônes}\begin{tabularx}{\linewidth}{@{}Yl@{}}\toprule
Reconnaître pyramides, cônes et solides usuels & \hyperref[exo-17-1]{17.1}, \hyperref[exo-17-5]{17.5}, \hyperref[exo-17-6]{17.6}\\[2mm]
Perspective et patrons & \hyperref[exo-17-9]{17.9}, \hyperref[exo-17-10]{17.10}, \hyperref[exo-17-11]{17.11}\\[2mm]
Volumes de la pyramide et du cône & \hyperref[exo-17-2]{17.2}, \hyperref[exo-17-3]{17.3}, \hyperref[exo-17-4]{17.4}, \hyperref[exo-17-7]{17.7}, \hyperref[exo-17-8]{17.8}, \hyperref[exo-17-12]{17.12}\\[2mm]
Unités, assemblages et problèmes & \hyperref[exo-17-13]{17.13}, \hyperref[exo-17-14]{17.14}, \hyperref[exo-17-15]{17.15}\\[2mm]
\bottomrule\end{tabularx}\par\vspace{4mm}
\otIntertitre{18. Algorithmique : conditions et variable}\begin{tabularx}{\linewidth}{@{}Yl@{}}\toprule
Variable, séquences et entrées & \hyperref[exo-18-1]{18.1}, \hyperref[exo-18-2]{18.2}, \hyperref[exo-18-5]{18.5}, \hyperref[exo-18-13]{18.13}\\[2mm]
Conditions et cas limites & \hyperref[exo-18-4]{18.4}, \hyperref[exo-18-6]{18.6}, \hyperref[exo-18-9]{18.9}, \hyperref[exo-18-15]{18.15}\\[2mm]
Boucles et constructions géométriques & \hyperref[exo-18-3]{18.3}, \hyperref[exo-18-7]{18.7}, \hyperref[exo-18-8]{18.8}, \hyperref[exo-18-11]{18.11}\\[2mm]
Événements, interactions et mise au point & \hyperref[exo-18-10]{18.10}, \hyperref[exo-18-12]{18.12}, \hyperref[exo-18-14]{18.14}\\[2mm]
\bottomrule\end{tabularx}\par\vspace{4mm}
\clearpage
\clearpage\phantomsection\label{matrice-objectifs}\pdfbookmark[1]{Objectifs officiels et chapitres}{bm-matrice-objectifs}\otPageTitre{Objectifs officiels et chapitres}{Annexe 2 du programme 2026 · rubriques de Quatrième}\fontsize{9.5}{12.5}\selectfont
\otIntertitre{Nombres et calculs}\begin{tabularx}{\linewidth}{@{}Yr@{}}\toprule
Multiplier deux relatifs, d’abord avec un seul facteur négatif, puis deux, grâce à la distributivité & \hyperref[chap-1]{ch. 1}\\[1.2mm]
Diviser deux relatifs ; calculer un enchaînement d’opérations & \hyperref[chap-1]{ch. 1}\\[1.2mm]
Utiliser le vocabulaire (somme, quotient\ldots) d’un programme de calcul et inversement & \hyperref[chap-1]{ch. 1}, \hyperref[chap-9]{ch. 9}\\[1.2mm]
Simplifier une fraction ; définir un nombre rationnel et exprimer son opposé & \hyperref[chap-2]{ch. 2}\\[1.2mm]
Produit de rationnels ; fraction d’une fraction, d’un nombre, d’une quantité & \hyperref[chap-2]{ch. 2}\\[1.2mm]
Inverse d’un nombre et d’une fraction ; division de fractions ; expressions et problèmes & \hyperref[chap-2]{ch. 2}\\[1.2mm]
Puissances d’exposant positif ; produits de puissances d’un même nombre ou de même exposant ; problèmes & \hyperref[chap-3]{ch. 3}\\[1.2mm]
Définition de la racine carrée ; encadrement par deux entiers consécutifs & \hyperref[chap-4]{ch. 4}\\[1.2mm]
Produire des formules et tester leur vraisemblance ; démontrer par le calcul algébrique & \hyperref[chap-5]{ch. 5}, \hyperref[chap-9]{ch. 9}\\[1.2mm]
Distributivité simple pour développer et factoriser & \hyperref[chap-5]{ch. 5}\\[1.2mm]
Résoudre $ax+b=c$ ; mettre en équation avec $ax+b=cx+d$ ; conjecturer avec un tableur ou un algorithme & \hyperref[chap-6]{ch. 6}\\[1.2mm]
\bottomrule\end{tabularx}\par\vspace{3mm}
\otIntertitre{Proportionnalité, fonctions}\begin{tabularx}{\linewidth}{@{}Yr@{}}\toprule
Grandeurs quotients, avec ou sans unités & \hyperref[chap-7]{ch. 7}, \hyperref[chap-8]{ch. 8}\\[1.2mm]
Comparer par un rapport ou un ratio ; égalités de rapports ; quatrième proportionnelle & \hyperref[chap-7]{ch. 7}\\[1.2mm]
Calculer avec des pourcentages ; coefficient multiplicateur & \hyperref[chap-8]{ch. 8}\\[1.2mm]
Résoudre des problèmes de partage proportionnel & \hyperref[chap-7]{ch. 7}\\[1.2mm]
Appliquer un programme de calcul à un nombre puis à une variable ; le remonter & \hyperref[chap-9]{ch. 9}\\[1.2mm]
Formule littérale, graphique et dépendance d’une grandeur en fonction d’une autre & \hyperref[chap-9]{ch. 9}\\[1.2mm]
\bottomrule\end{tabularx}\par\vspace{3mm}
\clearpage\otPageTitre{Objectifs officiels et chapitres}{Annexe 2 du programme 2026 · rubriques de Quatrième}\fontsize{9.5}{12.5}\selectfont
\otIntertitre{Organisation et gestion de données, probabilités}\begin{tabularx}{\linewidth}{@{}Yr@{}}\toprule
Moyenne pondérée d’une série de petit effectif & \hyperref[chap-10]{ch. 10}\\[1.2mm]
Médiane et étendue ; effet d’une valeur extrême ; problèmes et comparaison de séries & \hyperref[chap-10]{ch. 10}\\[1.2mm]
Utiliser le tableur pour la moyenne, la médiane et l’étendue & \hyperref[chap-10]{ch. 10}\\[1.2mm]
Vocabulaire et notations ensemblistes ; complémentaire, réunion, intersection, ensemble vide & \hyperref[chap-11]{ch. 11}\\[1.2mm]
Probabilité d’un événement et de l’événement contraire ; expériences à deux épreuves & \hyperref[chap-11]{ch. 11}\\[1.2mm]
Distributions fréquentielle et probabiliste ; fluctuation des fréquences & \hyperref[chap-11]{ch. 11}\\[1.2mm]
\bottomrule\end{tabularx}\par\vspace{3mm}
\otIntertitre{Espace et géométrie}\begin{tabularx}{\linewidth}{@{}Yr@{}}\toprule
Droites remarquables (automatismes) ; réciproque, contraposée et raisonnement par l’absurde & \hyperref[chap-12]{ch. 12}, \hyperref[chap-15]{ch. 15}\\[1.2mm]
Effet d’une translation ; lien avec les parallélogrammes et les angles ; conservations & \hyperref[chap-13]{ch. 13}\\[1.2mm]
Symétrie axiale, demi-tour, parallélogrammes particuliers (automatismes) & \hyperref[chap-13]{ch. 13}\\[1.2mm]
Les trois théorèmes relatifs à la droite des milieux & \hyperref[chap-14]{ch. 14}\\[1.2mm]
Théorème de Pythagore, réciproque et contraposée & \hyperref[chap-15]{ch. 15}\\[1.2mm]
Triangle rectangle et cercle circonscrit ; centre ; rectangles sans équerre & \hyperref[chap-16]{ch. 16}\\[1.2mm]
Pyramide et cône de révolution : reconnaître, représenter, calculer le volume & \hyperref[chap-17]{ch. 17}\\[1.2mm]
Repérage sur une droite graduée et dans un repère (automatismes) & \hyperref[chap-1]{ch. 1}, \hyperref[chap-9]{ch. 9}, \hyperref[chap-13]{ch. 13}\\[1.2mm]
\bottomrule\end{tabularx}\par\vspace{3mm}
\otIntertitre{Pensée informatique}\begin{tabularx}{\linewidth}{@{}Yr@{}}\toprule
Représenter des conditions simples ; écrire des instructions conditionnelles & \hyperref[chap-18]{ch. 18}\\[1.2mm]
Manipuler une variable ; écrire un programme simple ; modifier un programme donné & \hyperref[chap-18]{ch. 18}\\[1.2mm]
\bottomrule\end{tabularx}\par\vspace{3mm}
\clearpage
